% ============================================================================
% DISCUSSION --- four paragraphs planned:
%   1. what the paper establishes      (125-145 w)  [WRITTEN]
%   2. what is actually new            (105-120 w)  [WRITTEN]
%   3. limitations and boundaries      (145-165 w)  [WRITTEN]
%   4. outlook / broader implication   ( 80-100 w)  [WRITTEN]  -- Discussion complete.
% No separate Conclusion subsection.
% ============================================================================
\section{Discussion and limitations}
\label{sec:discussion}

Taken together, our results support a view of molecular design as control of a learned executable process rather than endpoint generation followed by scoring or repair. \method{} separates this process into three distinct objects: an executor that determines which molecular transitions are possible, a goal-independent reference law that learns which legal transitions are plausible, and a finite-horizon controller that favors successors according to what remains reachable under the current objective. The experiments show why this separation matters: the learned reference law captures state-dependent preferences beyond executable support alone, while future-aware control can improve molecular decisions without retraining that law. Because every committed state is itself a molecule, the same process can also be recontrolled from states already reached and restricted when admissibility must hold along the trajectory. The contribution is therefore not validity, molecular editing, or the $h$-transform in isolation, but a common molecular transition process on which plausibility and purpose remain separate.

This perspective also clarifies where \method{} differs from the ingredients on which it builds. Variable-size molecular generation and graph editing already provide mechanisms for changing molecular structure, while generator matching provides a general principle for learning stochastic processes and $h$-transforms provide a classical mechanism for conditioning them. \method{}'s distinction is in making these pieces act on the same molecular object: an explicit, state-dependent rewrite system defines the executable support, generator matching learns a goal-independent probability law over its canonical molecular successors, and finite-horizon control subsequently redirects that frozen law. This separation matters operationally because changing an objective need not redefine how molecules are represented, which transitions are legal, or which transitions the reference process has learned to prefer. Instead, task-specific information enters only when choosing among futures already supported by the learned process.

The same decomposition also makes the limitations of \method{} explicit. Control can redistribute probability only over transitions supported by the underlying molecular process; it cannot recover a useful continuation that is absent from the executable fiber or assigned negligible reference probability. Moreover, the exact finite-horizon guarantees apply when the backward value is exact, whereas molecular-scale control uses an amortized approximation $h_\phi$; empirical gains therefore demonstrate useful future-aware steering, not exact conditioning at scale. Property-directed experiments introduce a second dependency: both local and future-aware controllers inherit errors and extrapolation from the property models used to define molecular desirability. Finally, the executable state space is deliberately narrower than chemistry itself. Connectedness and valence admissibility in a 2D molecular graph do not establish stereochemical correctness, conformational plausibility, synthesizability, stability, or biological activity. These quantities require richer state representations or independent downstream evaluation rather than following from the closure of the rewrite system.

The natural next step is therefore not simply to make the controller stronger, but to expand the molecular process on which control operates. Richer executable rewrite languages could incorporate additional stereochemical, conformational, or reaction-aware state, while improved reference models could place useful probability mass on molecular transitions that are presently difficult to reach. The same separation between a goal-independent process and task-specific future value also makes experimental feedback conceptually straightforward: new property measurements can change what futures are preferred without requiring the molecular transition model to be redefined each time. More broadly, \method{} suggests a generative-modeling principle beyond the particular chemistry studied here: when a domain admits meaningful executable transformations between valid states, those transformations can serve as the native event language of a learned stochastic process, with generation and control operating on the same evolving object.
